\section{Context}

\subsection{Curiosity Origin}
As frequent travelers, we are curious about urban mobility and how people navigate large cities.
We became particularly interested in understanding mobility patterns in the context of sustainable transportation and green alternatives, recognizing their growing importance in urban planning and environmental goals.
Specifically, we wanted to explore the underlying choices people make when selecting transportation modes and the factors that shape those decisions.

This curiosity led us to explore data sources that could provide meaningful insights into urban movement.
Among them, New York City's bike rental data stood out as a rich source of information on how people travel across the city using a sustainable mode of transportation.
When combined with other datasets, it can offer a clearer picture of the city's daily dynamics.

We also learned that this bike rental data is provided by a private company (\texttt{Citi Bike}) that operates the bike-sharing service in New York City.
For this reason, making our analysis useful and relevant to the company is an effective way to increase its practical impact.

\subsection{Project Circumstance and Purpose}
Being one of the largest bike-sharing programs in the world, \texttt{Citi Bike} currently collects vast amounts of operational data. 
If effectively analyzed and visualized, this data could provide valuable insights to optimize their services and enhance user experience.

Our visualization tool directly addresses these gaps by making complex mobility data accessible and actionable.

The primary objective is to transform raw \texttt{Citi Bike} operational data into clear, actionable insights for stakeholders.

\subsubsection{Citi Bike Operations Team}
Specifically, we aim to reveal:
\begin{enumerate}
  \item \textbf{Temporal patterns}: peak usage hours, seasonal trends, and demand cyclicity.
  \item \textbf{Spatial hotspots}: high-demand stations, commute corridors, and underutilized infrastructure.
  \item \textbf{User segmentation}: behavioral differences between members and casual riders.
  \item \textbf{Rideable type preferences}: trends in rideable type usage (classic vs.\@ electric) across different contexts.
\end{enumerate}
The first two insights can help the operations team optimize station placement and bike rebalancing.
Meanwhile, the last two can inform targeted marketing strategies and service improvements to better meet user needs.

This information reflects the most relevant angles we aim to explore, as it directly informs operational decisions and strategic planning for the company.

\subsubsection{City Planners and General Public}
Beyond the operations team, our visualization can also benefit city planners by providing insights into commute 
flows and modal patterns, which can inform infrastructure investment and urban planning decisions.

For the general public, it can raise awareness of urban mobility and sustainable transportation alternatives,
fostering a more informed and engaged community.

\subsection{Intended Message and Driving Questions}
Beyond exposing the data for free exploration, our visualization argues a specific thesis:
\begin{quote}
\textit{Citi Bike has grown into a year-round engine of low-carbon mobility in New York City, and the size of that climate dividend depends on who is riding and on the weather.}
\end{quote}

This message is supported by three driving questions, ordered so that the narrative builds from scale, to nuance, to payoff:
\begin{enumerate}
  \item \textbf{Adoption: is bike usage growing?}
  We trace ridership since 2013 to show whether \texttt{Citi Bike} is becoming a more established part of the city's mobility, and whether the recent growth is increasingly electric.
  We report raw ride volumes while being explicit that part of any rise reflects fleet and station expansion, not demand alone.
  \item \textbf{Two riderships: who keeps riding when the weather turns?}
  We contrast members (largely commuters) with casual riders (largely leisure and tourism) to test the hypothesis that commuter demand is weather-resilient while casual demand is fair-weather, varying sharply with temperature, precipitation, and season.
  \item \textbf{Green dividend: how much car travel does this displace?}
  As the headline of the project, we translate ridership into an estimate of car-kilometres avoided and the associated \(\text{CO}_2\) savings, making the sustainability argument concrete.
  Because this estimate depends on an assumed car-substitution rate, we communicate it transparently as a range rather than a single figure.
\end{enumerate}

These questions nest rather than compete: the green dividend (Q3) scales with the volume established by adoption (Q1) and is modulated by the seasonal and user-type split revealed in (Q2).
This thesis-driven framing guides our default views, while the underlying exploratory interface still lets users investigate the data on their own terms.
